% 正方体 11 种展开图的 3 个代表：1-4-1 一字、T 字（1-4-1 变形）、Z 字（2-2-2）
\documentclass[tikz,border=4pt]{standalone}
\usepackage{ctex}
\usepackage{amsmath}
\usetikzlibrary{calc}
\begin{document}
\begin{tikzpicture}[thick, scale=0.5,
    sq/.style={draw, minimum size=1cm, inner sep=0}]
  \newcommand{\sq}[2]{\draw (#1,#2) rectangle ++(1,1);}

  % --- (a) 1-4-1 一字型：中间一排 4 个，左端上面 1 个、右端下面 1 个 ---
  \begin{scope}[xshift=0cm, yshift=0cm]
    \sq{0}{0}\sq{1}{0}\sq{2}{0}\sq{3}{0}
    \sq{0}{1}
    \sq{3}{-1}
    \node[font=\footnotesize] at (2,-2) {（a）1-4-1 型};
  \end{scope}

  % --- (b) T 型（1-4-1 的另一种变体）：中间一排 4，上方中部 1，下方中部 1 ---
  \begin{scope}[xshift=7cm, yshift=0cm]
    \sq{0}{0}\sq{1}{0}\sq{2}{0}\sq{3}{0}
    \sq{1}{1}
    \sq{2}{-1}
    \node[font=\footnotesize] at (2,-2) {（b）1-4-1 变形};
  \end{scope}

  % --- (c) Z 字型 2-2-2 ---
  \begin{scope}[xshift=14cm, yshift=0cm]
    \sq{0}{0}\sq{1}{0}
    \sq{1}{1}\sq{2}{1}
    \sq{2}{2}\sq{3}{2}
    \node[font=\footnotesize] at (2,-1) {（c）2-2-2 Z 型};
  \end{scope}
\end{tikzpicture}
\end{document}
